\section{A joint local limit with both collision endpoints}
\label{sec:collision-mixing-local-limit}
Compact-frequency estimates suffice for local intervals if the passage
to an unsmoothed interval preserves order. This section makes that
passage before inserting the exact flight-age overlap. In particular
we shall not substitute a polynomially increasing frequency cutoff
into an estimate whose constants were proved only on compact sets.

For $k\in\Z^2$, $t\in\R$ write
\[
 \zeta_{m,R}(k,t)=
 \left(\frac{k}{\sqrt m},\frac{t-m\bar\tau_R}{\sqrt m}\right),
 \qquad S_{m,R}=\sum_{j=0}^{m-1}\tau_R\circ T_R^j.
\]
All compact central sets below are fixed independently of $m$ and $R$.
The density $g_\Sigma$ is the normalized centered Gaussian in the
integer cell coordinates and the time coordinate.

\subsection{A fixed Fourier support and an order-preserving limit}
\begin{lemma}[Endpoint insertions with a band-limited time test]
\label{lem:bandlimited-collision-local}
Let $a,d$ be smooth collision functions. Suppose $q\in L^1(\R)$
has integrable Fourier transform of compact support, with the convention
$\widehat q(b)=\int e^{ibs}q(s)\,ds$. Then
\begin{align}
 &m^{3/2}\int a(x)d(T_R^m x)
   \one_{\{K^{\rm c}_{m,R}(x)=k\}}q(S_{m,R}(x)-t)\,d\nu(x)
 \notag\\
 &\qquad =g_{\Sigma_R}(\zeta_{m,R}(k,t))
             \left(\int a\,d\nu\right)
             \left(\int d\,d\nu\right)\int q(s)\,ds+o(1).
 \label{eq:bandlimited-collision-local}
\end{align}
The error is uniform in $R,k,t$ with $\zeta_{m,R}(k,t)$ in a fixed
compact set. For fixed $a,d,q$ the left side is bounded uniformly
in all $R,k,t,m\ge1$.
\end{lemma}
\begin{proof}
Torus orthogonality and Fourier inversion express the integral as
\begin{equation}\label{eq:bandlimited-collision-pairing}
 \frac1{(2\pi)^3}\int_{[-\pi,\pi]^2}\int_\R
 e^{-i[u\cdot k+b(t-m\bar\tau_R)]}\widehat q(-b)
 \ell\bigl(M_d Q_{R,(u,b)}^m(a\nu)\bigr)\,db\,du.
\end{equation}
The sign $-b$ in $\widehat q$ results from evaluating $q(S_m-t)$.
The equality follows directly from
\eqref{eq:exact-collision-endpoint-pairing}, without a change of the
endpoint functions or of the record.

The time frequency belongs to a fixed bounded interval. Outside a
fixed neighborhood of zero,
\eqref{eq:full-compact-frequency-gap} makes the integral exponentially
small, even after multiplication by $m^{3/2}$. Inside that neighborhood
use \eqref{eq:raw-collision-splitting} and put $v=\sqrt m(u,b)$.
The principal amplitude tends to $\nu(a)\nu(d)$; its value at zero
follows from $\Pi_R(0)=\nu\ell$. Also
$\widehat q(-v_3/\sqrt m)\to\int q$, uniformly on compact $v$ sets.
The absolute bound $C e^{-c|v|^2}$ for the principal term, and the
exponentially small complementary term on its fixed physical domain,
justify dominated convergence. The uniform expansion
\eqref{eq:raw-collision-expansion} identifies the limit as the Gaussian
Fourier inverse. The oscillatory factor has modulus one; convergence
is uniform on central compact sets, including along arbitrary
converging parameter sequences. The same absolute estimates bound the
normalized inverse for all $k,t$. Enlarge the constant for the
finitely many small values of $m$.
\end{proof}

\begin{lemma}[Fourier envelopes for an interval]
\label{lem:collision-interval-envelopes}
For a bounded interval $J$ and each $B>0$ there are real integrable
functions $q_{J,B}^-,q_{J,B}^+$ whose Fourier transforms are continuous
and supported in $[-B,B]$, such that
\begin{equation}\label{eq:collision-interval-envelopes}
 q_{J,B}^-\le\one_J\le q_{J,B}^+,
 \qquad \int q_{J,B}^{\pm}=|J|\pm\frac{2\pi}{B}.
\end{equation}
Endpoint conventions do not affect the inequalities if the lower
bound is taken for the open interval and the upper bound for its closure.
\end{lemma}
\begin{proof}
Use the Beurling--Selberg majorant and minorant for the interval;
we recall the normalization from \cite[Introduction]{CVExt}.
The Beurling functions $H,K$ satisfy
$|\operatorname{sgn}x-H(x)|\le K(x)$, where
$K(x)=(\sin(\pi x)/(\pi x))^2$, $\int K=1$, and $H$ is odd.
Their exponential type is $2\pi$. For $J$ with endpoints $c<d$ and
$\delta=B/(2\pi)$ take
\[
 q_{J,B}^{\pm}(s)=
 \tfrac12\{H(\delta(s-c))-H(\delta(s-d))\}
 \mathbin{\pm}\tfrac12\{K(\delta(s-c))+K(\delta(s-d))\}.
\]
The sign inequality gives the pointwise bounds. At an endpoint it gives
an upper value at least one and a lower value at most zero, so open,
closed and half-open intervals are all accommodated. The symmetric
integral of the difference of the two $H$ terms is $d-c$; the two
$K$ terms contribute $1/\delta$. The standard exponential-type
Fourier support statement gives $[-B,B]$ in the convention adopted
here. The functions are integrable; hence their transforms are
continuous, and compact support also makes those transforms integrable.
\end{proof}

\begin{theorem}[The mixed collision local limit with two endpoints]
\label{thm:two-endpoint-collision-LLT}
For continuous collision functions $a,d$ and any fixed bounded interval
$J$,
\begin{align}
 &m^{3/2}\int a(x)d(T_R^m x)
 \one_{\{K^{\rm c}_{m,R}=k,\ S_{m,R}-t\in J\}}\,d\nu(x)
 \notag\\
 &\qquad =|J|\,g_{\Sigma_R}(\zeta_{m,R}(k,t))\nu(a)\nu(d)+o(1),
 \label{eq:two-endpoint-collision-LLT}
\end{align}
uniformly in $R$ and on central compact sets. For each fixed bounded
interval $J$, independently of $k,t,R,m\ge1$,
\begin{equation}\label{eq:collision-local-upper-bound}
 m^{3/2}\nu\{K^{\rm c}_{m,R}=k,\ S_{m,R}-t\in J\}\le C_J.
\end{equation}
\end{theorem}
\begin{proof}
First let $a,d$ be nonnegative and smooth. Multiply the inequalities
\eqref{eq:collision-interval-envelopes} by the nonnegative source
weight $a(d\circ T_R^m)\one_{\{K_m^{\rm c}=k\}}$. For each fixed
$B$, Lemma~\ref{lem:bandlimited-collision-local} evaluates both bounds.
The difference of their limiting normalized values is
$4\pi B^{-1}g_{\Sigma_R}(\zeta)\nu(a)\nu(d)$, uniformly bounded
on the parameter interval. First take $m\to\infty$, then $B\to\infty$.
This proves the interval formula. In particular no growing-cutoff
spectral bound is used. The upper envelope for one fixed $B$, with
$a=d=1$, gives \eqref{eq:collision-local-upper-bound}.

Approximate nonnegative continuous $a,d$ uniformly by nonnegative
smooth functions on the compactified collision cylinder. The change
in the normalized integral is bounded by a constant times the
approximation error, by \eqref{eq:collision-local-upper-bound}.
This proves the result for continuous functions. Adding constants and
separating real and imaginary parts then gives the stated complex
version. Uniform convergence on central compact sets follows from the
uniform envelopes and the uniform fixed-band estimate.
\end{proof}

\subsection{Testing an endpoint-dependent overlap}
Let $\overline M$ be the compact collision cylinder, including the
null grazing boundary. Define positive, locally finite measures on
$\overline M\times\R\times\overline M$ by
\begin{equation}\label{eq:collision-local-source-measure}
 \rho_{m,R}^{k,t}
 =m^{3/2}\bigl(x,S_{m,R}(x)-t,T_R^m x\bigr)_*
                     (\one_{\{K^{\rm c}_{m,R}=k\}}\nu).
\end{equation}
\begin{corollary}[Local endpoint measures]
\label{cor:local-endpoint-measures}
If $R_m\to R_0$ and $\zeta_{m,R_m}(k_m,t_m)\to\zeta_0$, then
\begin{equation}\label{eq:collision-local-source-limit}
 \rho_{m,R_m}^{k_m,t_m}\longrightarrow
       g_{\Sigma_{R_0}}(\zeta_0)\,\nu(dx)\,ds\,\nu(dy)
\end{equation}
vaguely. This convergence also evaluates bounded, compactly
$s$-supported test functions continuous outside a set of zero measure
for the measure on the right.

It evaluates a varying test $F_{R_m}$ as well, provided there are open
neighborhoods $O_\eta$ whose closures have limiting measure tending
to zero as $\eta\downarrow0$ and whose boundaries have limiting
measure zero, such that $F_{R_m}\to F_{R_0}$ uniformly outside
$O_\eta$. The tests must have a common bound and a common compact
$s$ support. A small measure for the open set alone is not sufficient.
\end{corollary}
\begin{proof}
Products of continuous endpoint functions and interval indicators are
evaluated by Theorem~\ref{thm:two-endpoint-collision-LLT}. Their linear
combinations approximate continuous compactly supported tests, using
uniform continuity and a finite interval partition in $s$. The bound
\eqref{eq:collision-local-upper-bound} supplies uniform local masses.
This proves vague convergence. The limiting measure assigns zero mass
to every fixed time-slice boundary. Continuous upper and lower
approximations, or the Portmanteau theorem on a compact time slab,
then give the assertion about a fixed almost-everywhere continuous
test.

For the varying assertion use the exceptional neighborhoods with
small closed measure and null boundary required in the statement.
Portmanteau bounds the
limsup of the source measure of that closure by its limiting measure.
The common supremum bound controls both tests there. On its complement
uniform convergence controls their difference, by the local mass bound.
Take $m\to\infty$ and then $\eta\downarrow0$. Pointwise convergence
of the tests alone would not justify this step.
\end{proof}
